% !TeX encoding = UTF-8
% Reviewed conversion; see reports/revision-log.docx and bibliography.json.
\SourceChapter{3}{فصل سوم: روش پیشنهادی و معماری سیستم}
% SOURCE Introduction_Chapter3.txt:1-1 [chapter-title-in-wrapper]


% SOURCE Introduction_Chapter3.txt:2-2 [spacing]


% SOURCE Introduction_Chapter3.txt:3-3 [heading]
\SourceHeading{1}{۳{-}۱. مقدمه فصل سوم}

% SOURCE Introduction_Chapter3.txt:4-4 [paragraph]
در پاسخ به نیازهای مبرم امنیت، محرمانگی و صیانت از حریم خصوصی بیمار در سامانه اینترنت اشیای پزشکی \lr{(IoMT)} و چارچوب‌های پزشکی از راه دور \lr{(Telemedicine)}، در این فصل معماری جامع سیستم امنیتی هیبریدی پیشنهادی ارائه می‌گردد. این چارچوب چهارلایه جهت پوشش متوازن، رفع سه چالش بنیادین در امنیت دادگان سلامت دیجیتال و پاسخ‌گویی روش‌شناختی به فرضیات پژوهش و شکاف‌های استخراج‌شده در ادبیات تحقیق طراحى شده است \ThesisCite{1}, \ThesisCite{2}, \ThesisCite{8}:\par

% SOURCE Introduction_Chapter3.txt:5-5 [spacing]


% SOURCE Introduction_Chapter3.txt:6-6 [paragraph]
۱. تولید کلیدهای رمزی کاملاً پویا و وابسته به محتوا: با تلفیق شبکه عصبی عمیق \lr{U{-}Net} جهت قطعه‌بندی معنایی ناحیه حساس تشخیصی \lr{(ROI)} و تابع هش یک‌طرفه \lr{SHA{-}256} جهت مقداردهی پویای پارامترهای رمزی، که منحصراً به تصاویر ورودی وابسته بوده و امکان خنثی‌سازی کامل حملات متن‌واضح برگزیده \lr{(CPA)} و متن‌واضح معلوم \lr{(KPA)} را فراهم می‌سازد (پاسخ به فرضیه اول و شکاف پژوهشی اول).\par

% SOURCE Introduction_Chapter3.txt:7-7 [spacing]


% SOURCE Introduction_Chapter3.txt:8-8 [paragraph]
۲. شکستن همبستگی مکانی پیکسل‌ها و دستیابی به آنتروپی آرمانی: با بهره‌گیری از یک سیستم غیرخطی ابرآشوبی ۵ بعدی جدید (دارای بیش از یک توان نمای لیاپانوف مثبت)، آشفته‌سازی زیگ‌زاگی آدرس پیکسل‌ها و انتشار پویای دو مرحله‌ای \lr{DNA} (شامل \lr{Dynamic DNA Flip} و \lr{Dynamic DNA XOR}) جهت تبدیل ماتریس تصویر به نویز کاملاً یکنواخت و ایجاد پیچیدگی غیرخطی فوق‌العاده بالا (پاسخ به فرضیه دوم و شکاف پژوهشی سوم).\par

% SOURCE Introduction_Chapter3.txt:9-9 [spacing]


% SOURCE Introduction_Chapter3.txt:10-10 [paragraph]
۳. پنهان‌نگاری فراداده‌های هویتی بیمار با حفظ کامل اصالت و کیفیت تشخیصی بالینی: با فشرده‌سازی بدون اتلاف \lr{zlib} و جای‌دهی آگاه به برجستگی بصری \lr{(Saliency Detection)} فراداده‌های بیمار در نواحی کم‌اهمیت پس‌زمینه تصویر جهت حفظ کامل دست‌نخورده بودن بافت‌های تشخیصی (با هدف عددی \lr{PSNR \SourceGlyph{003E} 50 dB} و \lr{SSIM} \ensuremath{\approx} \lr{1.0} در ناحیه \lr{ROI}) (پاسخ به فرضیه سوم و شکاف پژوهشی دوم).\par

% SOURCE Introduction_Chapter3.txt:11-11 [spacing]


% SOURCE Introduction_Chapter3.txt:12-12 [paragraph]
ساختار یکپارچه این ۴ لایه پردازشی و نحوه تعامل آن‌ها در شکل ۳{-}۱ به تصویر کشیده شده است.\par

% SOURCE Introduction_Chapter3.txt:13-13 [spacing]


% SOURCE Introduction_Chapter3.txt:14-14 [image]
\SourcePicture{chapter3_system_architecture.png}{شکل ۳{-}۱: نمودار معماری مفهومی یکپارچه ۴لایه سیستم امنیتی هیبریدی پیشنهادی}{[شکل ۳{-}۱: نمودار معماری مفهومی یکپارچه ۴لایه سیستم امنیتی هیبریدی پیشنهادی {-} \lr{chapter3\_system\_architecture.png}]}{0}

% SOURCE Introduction_Chapter3.txt:15-15 [spacing]


% SOURCE Introduction_Chapter3.txt:16-16 [paragraph]
ساختار این فصل به گونه‌ای تنظیم شده است که ابتدا در بخش ۳{-}۲ (معماری کلی و گردش داده‌ها)، مدل فرستنده{-}گیرنده و فلوچارت‌های تفکیکی همگام‌سازی ارائه می‌شوند. در بخش‌های ۳{-}۳ تا ۳{-}۶ (لایه‌های چهارگانه)، جزئیات ریاضی، معادلات دیفرانسیل ۵بعدی، روابط \lr{DNA} و پنهان‌نگاری به همراه شبه‌کدهای چهارگانه اجرایی تبیین می‌گردند. در نهایت، در بخش ۳{-}۷ (تحلیل پیچیدگی محاسباتی)، کارایی زمانی سیستم با نماد \lr{O}{-}بزرگ و الزامات پیاده‌سازی زمان‌واقعی \lr{(Real{-}time)} بر روی سخت‌افزارهای لبه \lr{IoMT} (مانند \lr{NVIDIA Jetson TX2}) بررسی می‌شود. نقشه راه مفهومی و سیر ساختاری فصل سوم در شکل ۳{-}۲ نشان داده شده است.\par

% SOURCE Introduction_Chapter3.txt:17-17 [spacing]


% SOURCE Introduction_Chapter3.txt:18-18 [image]
\SourcePicture{chapter3_roadmap.png}{شکل ۳{-}۲: نمودار نقشه راه مفهومی و ساختاری فصل سوم (از تبیین معماری تا پیاده‌سازی لبه)}{[شکل ۳{-}۲: نمودار نقشه راه مفهومی و ساختاری فصل سوم (از تبیین معماری تا پیاده‌سازی لبه) {-} \lr{chapter3\_roadmap.png}]}{0}
% SOURCE Architecture_Dataflow_Chapter3.txt:1-1 [heading]
\SourceHeading{1}{۳{-}۲. معماری کلی چارچوب و گردش داده‌ها در فرستنده و گیرنده}

% SOURCE Architecture_Dataflow_Chapter3.txt:2-2 [spacing]


% SOURCE Architecture_Dataflow_Chapter3.txt:3-3 [paragraph]
برای پوشش هم‌زمان محرمانگی تصویر تشخیصی، صیانت از فراداده‌های هویتی بیمار و مقاومت در برابر حملات سایبری بر بستر شبکه‌های ناامن اینترنت اشیای پزشکی \lr{(IoMT)} و سامانه‌های پزشکی از راه دور \lr{(Telemedicine)}، چارچوب پیشنهادی بر پایه یک ساختار فرستنده{-}گیرنده متناظر و همگام‌سازی‌شده طراحی گردیده است. در این معماری، تصویر واضح ورودی $I$ با ابعاد $W \times H$ و فراداده متنی بیمار $M_{\mathrm{patient}}$ (شامل اطلاعات \lr{EXIF}، پرونده بالینی و شناسه هویتی) در سمت فرستنده پردازش شده و تصویر پنهان‌نگاری‌شده رمزشده $I_{\mathrm{Stego}}$ جهت ارسال بر روی شبکه تولید می‌شود \ThesisCite{1}, \ThesisCite{2}, \ThesisCite{8}.\par

% SOURCE Architecture_Dataflow_Chapter3.txt:4-4 [spacing]


% SOURCE Architecture_Dataflow_Chapter3.txt:5-5 [heading]
\SourceHeading{2}{۳{-}۲{-}۱. لوله پردازشی در سمت فرستنده \lr{(Sender Pipeline)}}

% SOURCE Architecture_Dataflow_Chapter3.txt:6-6 [spacing]


% SOURCE Architecture_Dataflow_Chapter3.txt:7-7 [paragraph]
زنجیره عملیاتی در سمت فرستنده شامل چهار گام پیاپی و وابسته به داده است که به صورت دقیق زیر فرمول‌بندی و اجرا می‌گردند:\par

% SOURCE Architecture_Dataflow_Chapter3.txt:8-8 [spacing]


% SOURCE Architecture_Dataflow_Chapter3.txt:9-9 [paragraph]
۱. پیش‌پردازش و استخراج خودکار ناحیه حساس \lr{(ROI Segmentation)}:\par

% SOURCE Architecture_Dataflow_Chapter3.txt:10-10 [paragraph]
تصویر پزشکی واضح ورودی $I$ ابتدا وارد شبکه عصبی عمیق \lr{U{-}Net} آموزش‌دیده می‌گردد. این شبکه با اعمال لایه‌های انقباضی و انبساطی، ماسک دوتایی ناحیه حساس تشخیصی $M_{\mathrm{ROI}}$ را با دقت بالا استخراج می‌نماید. ناحیه حساس تشخیصی شامل بافت‌های حیاتی بالینی (مانند ضایعات ریوی، تومورهای مغزی یا عروق شبکیه) است که هرگونه اعوجاج در آن‌ها منجر به خطای تشخیص پزشکی خواهد شد \ThesisCite{1}, \ThesisCite{3}.\par

% SOURCE Architecture_Dataflow_Chapter3.txt:11-11 [spacing]


% SOURCE Architecture_Dataflow_Chapter3.txt:12-12 [paragraph]
۲. استخراج ویژگی‌های آماری و کلیدسازی پویا \lr{(Dynamic Key Generation)}:\par

% SOURCE Architecture_Dataflow_Chapter3.txt:13-13 [paragraph]
ویژگی‌های آماری نابرابر ناحیه حساس استخراج‌شده (شامل میانگین شدت روشنایی، واریانس، کجی و کشیدگی هیستوگرام) محاسبه شده و به عنوان ورودی به تابع هش یک‌طرفه \lr{SHA{-}256} اعمال می‌گردند. چکیده ۲۵۶ بیتی حاصله ($K_{\mathrm{Send}}$)، به عنوان کلید پویا و کاملاً وابسته به تصویر عمل می‌کند. این ۲۵۶ بیت بر اساس یک مکانیزم نگاشت خطی به ۵ عدد اعشاری با دقت مضاعف تبدیل می‌شوند که شرایط اولیه $(x_0,\allowbreak{} y_0,\allowbreak{} z_0,\allowbreak{} u_0,\allowbreak{} v_0)$ و پارامترهای کنترل سیستم ابرآشوبی ۵ بعدی را مقداردهی اولیه می‌نمایند \ThesisCite{1}, \ThesisCite{4}.\par

% SOURCE Architecture_Dataflow_Chapter3.txt:14-14 [spacing]


% SOURCE Architecture_Dataflow_Chapter3.txt:15-15 [paragraph]
۳. رمزنگاری دو مرحله‌ای تصویر اصلی \lr{(Two{-}Stage Encryption: Confusion + Diffusion)}:\par

% SOURCE Architecture_Dataflow_Chapter3.txt:16-16 [paragraph]
{-} فاز آشفته‌سازی \lr{(Confusion)}: موقعیت مکانی پیکسل‌های تصویر واضح ورودی با استفاده از ماتریس جایگشت پیمایش زیگ‌زاگی \lr{(Zig{-}zag Scrambling)} جابه‌جا گردیده و همبستگی ساختاری پیکسل‌های همجوار شکسته می‌شود.\par

% SOURCE Architecture_Dataflow_Chapter3.txt:17-17 [paragraph]
{-} فاز انتشار \lr{(Diffusion)}: توالی‌های عددی شبه‌تصادفی حاصل از حل معادلات سیستم ابرآشوبی ۵ بعدی به توالی‌های \lr{DNA} نگاشت می‌شوند. سپس، دو فاز انتشار شامل چرخش پویای قواعد \lr{DNA} \lr{(Dynamic DNA Flip)} و عملگر منطقی \lr{Dynamic DNA XOR} بر روی ضرایب اعمال گردیده و تصویر رمزشده اولیه $I_{\mathrm{Enc}}$ با آنتروپی آرمانی شکل می‌گیرد \ThesisCite{1}, \ThesisCite{6}.\par

% SOURCE Architecture_Dataflow_Chapter3.txt:18-18 [spacing]


% SOURCE Architecture_Dataflow_Chapter3.txt:19-19 [paragraph]
۴. پنهان‌نگاری آگاه به برجستگی بصری فراداده \lr{(Saliency{-}Aware Steganography)}:\par

% SOURCE Architecture_Dataflow_Chapter3.txt:20-20 [paragraph]
فراداده متنی بیمار $M_{\mathrm{patient}}$ ابتدا توسط الگوریتم فشرده‌سازی بدون اتلاف \lr{zlib} فشرده شده و رشته باینری فشرده $P_{\mathrm{zlib}}$ را تشکیل می‌دهد. سپس، بر پایه نقشه برجستگی بصری \lr{(Saliency Detection)}، پیکسل‌های تصویر به دو دسته نواحی پرتوجه تشخیصی \lr{(High Saliency / ROI)} و نواحی کم‌برجستگی پس‌زمینه \lr{(Low Saliency / Background)} تقسیم می‌شوند. رشته فشرده فراداده صرفاً در بیت‌های کم‌ارزش \lr{(LSB)} پیکسل‌های پس‌زمینه جای‌دهی می‌شود تا تصویر نهایی $I_{\mathrm{Stego}}$ بدون ایجاد کوچک‌ترین اعوجاج در بافت‌های \lr{ROI} آماده ارسال گردد \ThesisCite{8}, \ThesisCite{9}.\par

% SOURCE Architecture_Dataflow_Chapter3.txt:21-21 [spacing]


% SOURCE Architecture_Dataflow_Chapter3.txt:22-22 [image]
\SourcePicture{chapter3_sender_flowchart.png}{شکل ۳{-}۳: فلوچارت تفصیلی مراحل چهارگانه رمزنگاری و پنهان‌نگاری در سمت فرستنده}{[شکل ۳{-}۳: فلوچارت تفصیلی مراحل چهارگانه رمزنگاری و پنهان‌نگاری در سمت فرستنده {-} \lr{chapter3\_sender\_flowchart.png}]}{0}

% SOURCE Architecture_Dataflow_Chapter3.txt:23-23 [spacing]


% SOURCE Architecture_Dataflow_Chapter3.txt:24-24 [heading]
\SourceHeading{2}{۳{-}۲{-}۲. قضیه ریاضی هم‌ارزی و بازسازی پویای کلید رمزی در سمت گیرنده \lr{(Key Reversibility Theorem)}}

% SOURCE Architecture_Dataflow_Chapter3.txt:25-25 [spacing]


% SOURCE Architecture_Dataflow_Chapter3.txt:26-26 [paragraph]
یکی از چالش‌های بنیادی در سیستم‌های رمزنگاری مبتنی بر کلید پویا، نحوه ارسال یا بازسازی کلید رمزی در سمت گیرنده بدون ایجاد کانال جانبی \lr{(Side Channel)} و بدون افشای کلید بر روی شبکه است. در معماری پیشنهادی، این چالش بر اساس قضیه هم‌ارزی بازسازی کلید اثبات و حل گردیده است:\par

% SOURCE Architecture_Dataflow_Chapter3.txt:27-27 [spacing]


% SOURCE Architecture_Dataflow_Chapter3.txt:28-28 [paragraph]
از آنجایی که فاز پنهان‌نگاری فراداده صراحتاً محدود به نواحی کم‌برجستگی پس‌زمینه بوده و ناحیه تشخیصی \lr{ROI} ۱۰۰\SourceGlyph{066A} دست‌نخورده باقی می‌ماند ($\mathrm{ROI} \cap \mathrm{Payload} = \emptyset$)، پیکسل‌های ناحیه \lr{ROI} در تصویر استگوی دریافتی $I_{\mathrm{Stego}}$ دقیقاً برابر با پیکسل‌های \lr{ROI} در تصویر واضح اولیه $I$ می‌باشند. بنابراین، با اعمال شبکه \lr{U{-}Net} در سمت گیرنده، ماسک \lr{ROI} و ویژگی‌های آماری استخراج‌شده کاملاً متناظر با فرستنده بوده و چکیده ۲۵۶ بیتی \lr{SHA{-}256} و شرایط اولیه سیستم ۵ بعدی عینا و بدون کوچک‌ترین خطایی بازسازی می‌گردند ($K_{\mathrm{Rec}} \equiv K_{\text{Send}}$) \ThesisCite{1}, \ThesisCite{4}. با این حال، به دلیل تغییرات جزئی ایجادشده در لایه‌های کم‌ارزش پس‌زمینه تصویر استگو توسط فرآیند پنهان‌نگاری، بازسازی بیت‌به‌بیت و قطعی ناحیه حساس بالینی مستقیماً نیازمند دریافت یک پرونده کمکی احرازهویت‌شده موسوم به \lr{Recovery Sidecar} (\texttt{recovery.msr}) است که حاوی مقادیر اصلی بیت‌های جایگزین‌شده و تگ احرازهویت \lr{AES-GCM-256} بوده و در صورت صحت‌آزمایی، دسترسی به تصویر اصلی بدون هیچ‌گونه اتلاف اطلاعات بالینی ($\text{PSNR}_{\mathrm{ROI}} = \infty$) میسر می‌سازد \ThesisCite{2}.\par

% SOURCE Architecture_Dataflow_Chapter3.txt:29-29 [spacing]


% SOURCE Architecture_Dataflow_Chapter3.txt:30-30 [image]
\SourcePicture{eq3_0_key_reversibility.png}{رابطه ۳{-}۱: قضیه ریاضی هم‌ارزی و بازسازی پویای کلید رمزی در سمت گیرنده}{[رابطه ۳{-}۱: قضیه ریاضی هم‌ارزی و بازسازی پویای کلید رمزی در سمت گیرنده {-} \lr{eq3\_0\_key\_reversibility.png}]}{2}

% SOURCE Architecture_Dataflow_Chapter3.txt:31-31 [spacing]


% SOURCE Architecture_Dataflow_Chapter3.txt:32-32 [heading]
\SourceHeading{2}{۳{-}۲{-}۳. لوله پردازشی در سمت گیرنده \lr{(Receiver Pipeline)}}

% SOURCE Architecture_Dataflow_Chapter3.txt:33-33 [spacing]


% SOURCE Architecture_Dataflow_Chapter3.txt:34-34 [paragraph]
در سمت گیرنده، زنجیره معکوس بدون نیاز به تصویر اصلی اولیه و بدون نیاز به دریافت کلید از کانال مجزا، بر روی تصویر $I_{\mathrm{Stego}}$ به صورت زیر اجرا می‌گردد:\par

% SOURCE Architecture_Dataflow_Chapter3.txt:35-35 [spacing]


% SOURCE Architecture_Dataflow_Chapter3.txt:36-36 [paragraph]
۱. استخراج فراداده بیمار و بازسازی \lr{zlib}:\par

% SOURCE Architecture_Dataflow_Chapter3.txt:37-37 [paragraph]
بر پایه الگوریتم برجستگی بصری، پیکسل‌های پس‌زمینه شناسایی شده و بیت‌های \lr{LSB} استخراج می‌گردند. رشته باینری حاصله توسط الگوریتم \lr{zlib} واسازی \lr{(Decompress)} شده و متن پرونده بیمار $M_{\mathrm{patient}}$ با صحت ۱۰۰\SourceGlyph{066A} بازیابی می‌شود \ThesisCite{8}, \ThesisCite{9}.\par

% SOURCE Architecture_Dataflow_Chapter3.txt:38-38 [spacing]


% SOURCE Architecture_Dataflow_Chapter3.txt:39-39 [paragraph]
۲. استخراج \lr{ROI}، محاسبه کلید معتبر و تولید توالی‌های ۵ بعدی:\par

% SOURCE Architecture_Dataflow_Chapter3.txt:40-40 [paragraph]
شبکه عصبی \lr{U{-}Net} بر روی تصویر دریافتی اعمال شده و ناحیه دست‌نخورده \lr{ROI} استخراج می‌گردد. ویژگی‌های آماری \lr{ROI} محاسبه شده و کلید ۲۵۶ بیتی \lr{SHA{-}256} تولید می‌شود. شرایط اولیه $(x_0,\allowbreak{} y_0,\allowbreak{} z_0,\allowbreak{} u_0,\allowbreak{} v_0)$ بازسازی شده و توالی‌های ۵ بعدی با حل عددی \lr{RK4} مجدداً تولید می‌گردند \ThesisCite{1}.\par

% SOURCE Architecture_Dataflow_Chapter3.txt:41-41 [spacing]


% SOURCE Architecture_Dataflow_Chapter3.txt:42-42 [paragraph]
۳. انتشار معکوس \lr{DNA}، جایگشت معکوس زیگ‌زاگی و بازسازی ۱۰۰\SourceGlyph{066A} تصویر واضح:\par

% SOURCE Architecture_Dataflow_Chapter3.txt:43-43 [paragraph]
توالی‌های ۵ بعدی به قواعد \lr{DNA} نگاشت شده و عملگرهای معکوس \lr{Dynamic DNA XOR} و معکوس چرخش \lr{DNA} بر روی تصویر اعمال می‌شوند. در نهایت، با اجرای جایگشت معکوس زیگ‌زاگی \lr{(Inverse Zig{-}zag Scrambling)}، تصویر واضح اصلی $I$ به صورت کاملاً بدون اتلاف ($PSNR = \infty$ و $SSIM = 1.0$) بازیابی می‌گردد \ThesisCite{1}, \ThesisCite{6}.\par

% SOURCE Architecture_Dataflow_Chapter3.txt:44-44 [spacing]


% SOURCE Architecture_Dataflow_Chapter3.txt:45-45 [image]
\SourcePicture{chapter3_receiver_flowchart.png}{شکل ۳{-}۴: فلوچارت تفصیلی مراحل سه‌گانه بازسازی، استخراج فراداده و رمزگشایی در سمت گیرنده}{[شکل ۳{-}۴: فلوچارت تفصیلی مراحل سه‌گانه بازسازی، استخراج فراداده و رمزگشایی در سمت گیرنده {-} \lr{chapter3\_receiver\_flowchart.png}]}{0}
% SOURCE UNet_KeyGen_Chapter3.txt:1-1 [heading]
\SourceHeading{1}{۳{-}۳. ماژول قطعه‌بندی معنایی \lr{U{-}Net}، استخراج ویژگی و تولید کلید پویای \lr{SHA{-}256}}

% SOURCE UNet_KeyGen_Chapter3.txt:2-2 [spacing]


% SOURCE UNet_KeyGen_Chapter3.txt:3-3 [paragraph]
جهت شکستن وابستگی به کلیدهای ثابت و ارائه یک الگوریتم مقاوم در برابر حملات دیفرانسیلی، متن‌واضح برگزیده \lr{(CPA)} و متن‌واضح معلوم \lr{(KPA)}، سیستم امنیتی پیشنهادی از ویژگی‌های ساختاری و آماری خودِ تصویر ورودی برای تولید کلید رمزی پویا استفاده می‌کند. با این حال، استفاده مستقیم از تمام پیکسل‌های تصویر اصلی باعث می‌شود که سیستم در برابر کوچک‌ترین نویزهای محیطی یا تغییرات ناخواسته در پس‌زمینه نامقاوم گردد. برای حل این مشکل بنیادی، استخراج ویژگی‌ها صرفاً بر روی ناحیه حساس تشخیصی \lr{(ROI)} متمرکز گردیده است \ThesisCite{1}, \ThesisCite{2}, \ThesisCite{3}.\par

% SOURCE UNet_KeyGen_Chapter3.txt:4-4 [spacing]


% SOURCE UNet_KeyGen_Chapter3.txt:5-5 [paragraph]
در گام اول، تصویر پزشکی ورودی $I$ با ابعاد $W \times H$ به شبکه عصبی عمیق \lr{U{-}Net} آموزش‌دیده تغذیه می‌شود. شبکه \lr{U{-}Net} با استفاده از معماری انقباضی{-}انبساطی و اتصالات پرشی \lr{(Skip Connections)}، ماسک دوتایی ناحیه حساس تشخیصی $M_{\mathrm{ROI}}$ را با دقت بالا استخراج می‌نماید:\par

% SOURCE UNet_KeyGen_Chapter3.txt:6-6 [spacing]


% SOURCE UNet_KeyGen_Chapter3.txt:7-7 [paragraph]
\[M_{\mathrm{ROI}} = \mathrm{U\text{-}Net}(I),\allowbreak{} \quad M_{\mathrm{ROI}}(i,\allowbreak{}j) \in \{0,\allowbreak{} 1\}\]\par

% SOURCE UNet_KeyGen_Chapter3.txt:8-8 [spacing]


% SOURCE UNet_KeyGen_Chapter3.txt:9-9 [image]
\SourcePicture{chapter3_unet_roi_key_extraction.png}{شکل ۳{-}۵: نمودار بلاکی فرآیند قطعه‌بندی معنایی \lr{U{-}Net}، استخراج ویژگی‌های آماری \lr{ROI} و محاسبه کلید ۲۵۶ بیتی \lr{SHA{-}256}}{[شکل ۳{-}۵: نمودار بلاکی فرآیند قطعه‌بندی معنایی \lr{U{-}Net}، استخراج ویژگی‌های آماری \lr{ROI} و محاسبه کلید ۲۵۶ بیتی \lr{SHA{-}256 {-} chapter3\_unet\_roi\_key\_extraction.png}]}{0}

% SOURCE UNet_KeyGen_Chapter3.txt:10-10 [spacing]


% SOURCE UNet_KeyGen_Chapter3.txt:11-11 [paragraph]
در گام دوم، پیکسل‌های واقع در ناحیه $M_{\mathrm{ROI}}$ تفکیک شده و چهار گشتاور آماری کلیدی هیستوگرام روشنایی جهت توصیف منحصر‌به‌فرد تصویر استخراج می‌گردند:\par

% SOURCE UNet_KeyGen_Chapter3.txt:12-12 [paragraph]
۱. میانگین شدت روشنایی پیکسل‌های \lr{ROI} ($\mu_{\mathrm{ROI}}$): نشان‌دهنده سطح تباین متوسط بافت تشخیصی.\par

% SOURCE UNet_KeyGen_Chapter3.txt:13-13 [paragraph]
۲. واریانس و انحراف معیار شدت پیکسل‌ها ($\sigma_{\mathrm{ROI}}^2$): نشان‌دهنده پراکندگی مقادیر روشنایی در ناحیه آسیب‌دیده.\par

% SOURCE UNet_KeyGen_Chapter3.txt:14-14 [paragraph]
۳. گشتاور سوم یا کجی هیستوگرام ($S_{\mathrm{ROI}}$): سنجش عدم تقارن توزیع آماری پیکسل‌های \lr{ROI}.\par

% SOURCE UNet_KeyGen_Chapter3.txt:15-15 [paragraph]
۴. گشتاور چهارم یا کشیدگی هیستوگرام ($K_{\mathrm{ROI}}$): سنجش میزان تیز بودن قله هیستوگرام و حضور نویزهای نقطه‌ای.\par

% SOURCE UNet_KeyGen_Chapter3.txt:16-16 [spacing]


% SOURCE UNet_KeyGen_Chapter3.txt:17-17 [paragraph]
بردار ویژگی‌های آماری حاصل به همراه ابعاد تصویر به عنوان ورودی به تابع هش یک‌طرفه و امن \lr{SHA{-}256} تغذیه می‌شوند تا چکیده باینری ۲۵۶ بیتی $K_{\mathrm{hash}}$ تولید گردد. به منظور جلوگیری از خطاهای ناشی از گردکردن اعشاری در محیط‌های محاسباتی مختلف، مقادیر گشتاورهای آماری با دقت ممیز شناور ۶۴ بیتی و ۱۷ رقم بااهمیت \lr{(IEEE 754 double-precision)} و با استفاده از کاراکتر جداکننده خط عمودی (\texttt{|}) و درج صریح ابعاد تصویر سریال‌سازی می‌شوند:\par

% SOURCE UNet_KeyGen_Chapter3.txt:18-18 [spacing]


% SOURCE UNet_KeyGen_Chapter3.txt:19-19 [paragraph]
\[\mathcal{F}_{\mathrm{str}} = \mathrm{Str}(\mu) \parallel \texttt{|} \parallel \mathrm{Str}(\sigma^2) \parallel \texttt{|} \parallel \mathrm{Str}(\gamma) \parallel \texttt{|} \parallel \mathrm{Str}(\kappa) \parallel \texttt{|} \parallel \mathrm{Str}(W) \parallel \texttt{|} \parallel \mathrm{Str}(H)\]
\[K_{\mathrm{hash}} = \mathrm{SHA256}(\mathcal{F}_{\mathrm{str}})\]\par

% SOURCE UNet_KeyGen_Chapter3.txt:20-20 [spacing]


% SOURCE UNet_KeyGen_Chapter3.txt:21-21 [paragraph]
در گام سوم، کلید باینری ۲۵۶ بیتی $K_{\mathrm{hash}}$ به ۵ بلوک ۴۸ بیتی مجزا ($K_{b1},\allowbreak{} K_{b2},\allowbreak{} K_{b3},\allowbreak{} K_{b4},\allowbreak{} K_{b5}$) تقسیم شده و از طریق روابط ریاضی مشخص، به ۵ عدد اعشاری پیوسته در بازه متقارن $[-2.0, +2.0]$ نگاشت می‌شوند تا شرایط اولیه $(x_0,\allowbreak{} y_0,\allowbreak{} z_0,\allowbreak{} u_0,\allowbreak{} v_0)$ برای راه‌اندازی دستگاه معادلات دیفرانسیل ۵بعدی ابرآشوبی حاصل گردد. این نگاشت خطی به صورت $s_0 = \left( \frac{K_b \bmod 10^8}{10^8} \right) \times 4.0 - 2.0$ انجام شده و تضمین می‌کند که متغیرهای حالت بلافاصله در حوضه جاذب غریب \lr{(Strange Attractor)} سیستم قرار می‌گیرند. فرمول‌بندی کامل این فرآیند در رابطه ۳{-}۲ ارائه شده است.\par

% SOURCE UNet_KeyGen_Chapter3.txt:22-22 [spacing]


% SOURCE UNet_KeyGen_Chapter3.txt:23-23 [image]
\SourcePicture{eq3_1_sha256_key_gen.png}{رابطه ۳{-}۲: فرمول‌بندی ریاضی استخراج ویژگی‌های آماری \lr{ROI}، محاسبه چکیده \lr{SHA{-}256} و نرمال‌سازی شرایط اولیه سیستم ۵بعدی}{[رابطه ۳{-}۲: فرمول‌بندی ریاضی استخراج ویژگی‌های آماری \lr{ROI}، محاسبه چکیده \lr{SHA{-}256} و نرمال‌سازی شرایط اولیه سیستم ۵بعدی {-} \lr{eq3\_1\_sha256\_key\_gen.png}]}{2}

% SOURCE UNet_KeyGen_Chapter3.txt:24-24 [spacing]


% SOURCE UNet_KeyGen_Chapter3.txt:25-25 [paragraph]
الگوریتم گام‌به‌گام ماژول قطعه‌بندی \lr{U{-}Net} و استخراج کلید پویا در شبه‌کد ۳{-}۱ آورده شده است.\par

% SOURCE UNet_KeyGen_Chapter3.txt:26-26 [spacing]


% SOURCE UNet_KeyGen_Chapter3.txt:27-49 [pseudocode]
\Needspace{5\baselineskip}\noindent [شبه‌کد ۳{-}۱: قطعه‌بندی \lr{ROI} با \lr{U{-}Net}، استخراج ویژگی‌های آماری و کلیدسازی پویا با \lr{SHA{-}256}]\par\begin{sourcecode}
\begin{LTR}\ttfamily\fontsize{9}{13}\selectfont\raggedright
\CodeLine{0.0}{Input  : Medical Image I (W x H), Trained U{-}Net Model}
\CodeLine{0.0}{Output : Binary Mask M\_\allowbreak{}ROI, 256{-}bit Hash Key K\_\allowbreak{}hash, 5D Initial Values (x0, y0, z0, u0, v0)}
\CodeLine{0.0}{\strut}
\CodeLine{0.0}{1. M\_\allowbreak{}ROI = U\_\allowbreak{}Net\_\allowbreak{}Inference(I, Trained\_\allowbreak{}Model)}
\CodeLine{0.0}{2. Extract ROI Pixels: P\_\allowbreak{}ROI = \{ I(i,j) \SourceGlyph{007C} M\_\allowbreak{}ROI(i,j) == 1 \}}
\CodeLine{0.0}{3. Compute Statistical Moments:}
\CodeLine{1.5}{{-} mu\_\allowbreak{}ROI   = Mean(P\_\allowbreak{}ROI)}
\CodeLine{1.5}{{-} var\_\allowbreak{}ROI  = Variance(P\_\allowbreak{}ROI)}
\CodeLine{1.5}{{-} skew\_\allowbreak{}ROI = Skewness(P\_\allowbreak{}ROI)}
\CodeLine{1.5}{{-} kurt\_\allowbreak{}ROI = Kurtosis(P\_\allowbreak{}ROI)}
\CodeLine{0.0}{4. Construct Feature Stream: F\_\allowbreak{}str = String(mu\_\allowbreak{}ROI) \SourceGlyph{007C}\SourceGlyph{007C} String(var\_\allowbreak{}ROI) \SourceGlyph{007C}\SourceGlyph{007C} String(skew\_\allowbreak{}ROI) \SourceGlyph{007C}\SourceGlyph{007C} String(kurt\_\allowbreak{}ROI) \SourceGlyph{007C}\SourceGlyph{007C} String(W) \SourceGlyph{007C}\SourceGlyph{007C} String(H)}
\CodeLine{0.0}{5. K\_\allowbreak{}hash = SHA256(F\_\allowbreak{}str)  // 256{-}bit hexadecimal / binary stream}
\CodeLine{0.0}{6. Split K\_\allowbreak{}hash into 5 blocks of 48 bits: K\_\allowbreak{}b1, K\_\allowbreak{}b2, K\_\allowbreak{}b3, K\_\allowbreak{}b4, K\_\allowbreak{}b5}
\CodeLine{0.0}{7. Compute Initial Values:}
\CodeLine{1.5}{{-} x0 = (K\_\allowbreak{}b1 mod 10\SourceGlyph{005E}8) / 10\SourceGlyph{005E}8}
\CodeLine{1.5}{{-} y0 = (K\_\allowbreak{}b2 mod 10\SourceGlyph{005E}8) / 10\SourceGlyph{005E}8}
\CodeLine{1.5}{{-} z0 = (K\_\allowbreak{}b3 mod 10\SourceGlyph{005E}8) / 10\SourceGlyph{005E}8}
\CodeLine{1.5}{{-} u0 = (K\_\allowbreak{}b4 mod 10\SourceGlyph{005E}8) / 10\SourceGlyph{005E}8}
\CodeLine{1.5}{{-} v0 = (K\_\allowbreak{}b5 mod 10\SourceGlyph{005E}8) / 10\SourceGlyph{005E}8}
\CodeLine{0.0}{8. Return M\_\allowbreak{}ROI, K\_\allowbreak{}hash, (x0, y0, z0, u0, v0)}
\end{LTR}
\end{sourcecode}

% SOURCE UNet_KeyGen_Chapter3.txt:50-50 [spacing]

% SOURCE Hyperchaos_5D_Chapter3.txt:1-1 [heading]
\SourceHeading{1}{۳{-}۴. سیستم غیرخطی ابرآشوبی ۵ بعدی، حل عددی و تولید کلیدهای رمزی}

% SOURCE Hyperchaos_5D_Chapter3.txt:2-2 [spacing]


% SOURCE Hyperchaos_5D_Chapter3.txt:3-3 [paragraph]
برای تضمین آنتروپی آرمانی (\lr{H}\lr{(m)} \ensuremath{\approx} \lr{8.0})، شکستن کامل همبستگی مکانی پیکسل‌های مجاور، مقاومت در برابر حملات بازسازی فضای فاز \lr{(Phase Space Reconstruction Attack)} و ایجاد فضای کلید عظیم فراتر از \lr{2}\SourceGlyph{005E}\lr{520}، سیستم امنیتی پیشنهادی از یک سیستم غیرخطی ابرآشوبی ۵ بعدی پیوسته جدید استفاده می‌نماید \ThesisCite{1}. در مقایسه با نگاشت‌های آشوبی کم‌بعد سنتی (نظیر نگاشت لوجستیک ۱ بعدی یا نگاشت هنون ۲ بعدی) که به دلیل سادگی ساختار ریاضی در برابر حملات تحلیل تفاضلی و \lr{CPA} آسیب‌پذیر هستند، سیستم ۵ بعدی به دلیل دارا بودن بیش از یک توان نمای لیاپانوف مثبت \lr{(Multiple Positive Lyapunov Exponents: \ensuremath{\lambda}1, \ensuremath{\lambda}2 \SourceGlyph{003E} 0)}، پیچیدگی دینامیکی فوق‌العاده بالا و رفتار شبه‌تصادفی غیردوره‌ای، بالاترین سطح محرمانگی را برای تصاویر پزشکی فراهم می‌سازد \ThesisCite{1}, \ThesisCite{6}, \ThesisCite{7}.\par

% SOURCE Hyperchaos_5D_Chapter3.txt:4-4 [spacing]


% SOURCE Hyperchaos_5D_Chapter3.txt:5-5 [heading]
\SourceHeading{2}{۳{-}۴{-}۱. فرمول‌بندی ریاضی دستگاه معادلات دیفرانسیل سیستم ۵بعدی}

% SOURCE Hyperchaos_5D_Chapter3.txt:6-6 [paragraph]
دستگاه معادلات دیفرانسیل حالت پیوسته سیستم ۵ بعدی ابرآشوبی پیشنهادی، بر پایه مدل سوباترا و تانیکایسلَون (\cite{Subathra2025}) و شامل ۵ متغیر حالت $(x, y, z, u, v)$ و ۷ پارامتر کنترلی غیرخطی به صورت زیر فرمول‌بندی می‌گردد:\par

% SOURCE Hyperchaos_5D_Chapter3.txt:7-7 [spacing]


% SOURCE Hyperchaos_5D_Chapter3.txt:8-8 [image]
\SourcePicture{eq3_2_5d_hyperchaos_rk4.png}{رابطه ۳{-}۳: دستگاه معادلات دیفرانسیل پیوسته سیستم ۵ بعدی ابرآشوبی و الگوریتم حل عددی \lr{RK4}}{[رابطه ۳{-}۳: دستگاه معادلات دیفرانسیل پیوسته سیستم ۵ بعدی ابرآشوبی و الگوریتم حل عددی \lr{RK4 {-} eq3\_2\_5d\_hyperchaos\_rk4.png}]}{2}

% SOURCE Hyperchaos_5D_Chapter3.txt:9-9 [spacing]


% SOURCE Hyperchaos_5D_Chapter3.txt:10-10 [paragraph]
که در آن، شرایط اولیه سیستم \lr{(x0, y0, z0, u0, v0)} مستقیماً از خروجی نرمال‌سازی‌شده تابع هش \lr{SHA{-}256} (محاسبه‌شده از ویژگی‌های آماری \lr{ROI} در بخش ۳{-}۳) تغذیه می‌شوند. این امر وابستگی ۱۰۰ درصدی توالی‌های رمزی حاصله به تصویر واضح ورودی را تضمین می‌نماید.\par

% SOURCE Hyperchaos_5D_Chapter3.txt:11-11 [spacing]


% SOURCE Hyperchaos_5D_Chapter3.txt:12-12 [paragraph]
برای تجسم رفتار دینامیکی، پیچیدگی فضای فاز و اثبات وجود جاذبه غریب \lr{(Strange Attractor)}، تصوير پرتره فاز ۳ بعدی حاصل از تصویر کردن متغیرهای حالت \lr{(x, y, z)} در شکل ۳{-}۶ نشان داده شده است.\par

% SOURCE Hyperchaos_5D_Chapter3.txt:13-13 [spacing]


% SOURCE Hyperchaos_5D_Chapter3.txt:14-14 [image]
\SourcePicture{chapter3_5d_hyperchaos_phase_portrait.png}{شکل ۳{-}۶: پرتره فاز ۳ بعدی و نمایش جاذبه غریب سیستم ابرآشوبی ۵ بعدی در فضای حالت \lr{(x{-}y{-}z)}}{[شکل ۳{-}۶: پرتره فاز ۳ بعدی و نمایش جاذبه غریب سیستم ابرآشوبی ۵ بعدی در فضای حالت \lr{(x{-}y{-}z)} {-} \lr{chapter3\_5d\_hyperchaos\_phase\_portrait.png}]}{0}

% SOURCE Hyperchaos_5D_Chapter3.txt:15-15 [spacing]


% SOURCE Hyperchaos_5D_Chapter3.txt:16-16 [heading]
\SourceHeading{2}{۳{-}۴{-}۲. حل عددی دستگاه معادلات با روش رونگه{-}کوتا مرتبه ۴ \lr{(RK4)}}

% SOURCE Hyperchaos_5D_Chapter3.txt:17-17 [paragraph]
از آنجایی که معادلات دیفرانسیل سیستم ۵ بعدی پیوسته هستند، برای پیاده‌سازی دیجیتال و استخراج توالی‌های رمزی گسسته، دستگاه معادلات با استفاده از روش حل عددی رونگه{-}کوتا مرتبه ۴ \lr{(Runge{-}Kutta 4th Order {-} RK4)} با گام زمانی یکنواخت \lr{h = 0.001} حل می‌گردد.\par

% SOURCE Hyperchaos_5D_Chapter3.txt:18-18 [spacing]


% SOURCE Hyperchaos_5D_Chapter3.txt:19-19 [paragraph]
برای حذف اثرات گذرا \lr{(Transient Effect)} و اطمینان از قرارگیری کامل سیستم در فاز ابرآشوبی پایدار، \lr{N0 = 1000} نمونه اولیه تولیدشده از توالی حذف گردیده و مابقی نمونه‌ها به تعداد لازم برای ابعاد تصویر \lr{(M = W x H)} استخراج می‌شوند.\par

% SOURCE Hyperchaos_5D_Chapter3.txt:20-20 [spacing]


% SOURCE Hyperchaos_5D_Chapter3.txt:21-21 [heading]
\SourceHeading{2}{۳{-}۴{-}۳. گسسته‌سازی و استخراج کلیدهای رمزی ۸ بیتی (\lr{K1} تا \lr{K5})}

% SOURCE Hyperchaos_5D_Chapter3.txt:22-22 [paragraph]
توالی‌های پیوسته اعشاری حاصل از حل عددی \lr{(x\_i, y\_i, z\_i, u\_i, v\_i)} از طریق نگاشت اعشاری به اعداد صحیح ۸ بیتی در بازه \lr{[0, 255]} تبدیل می‌شوند تا ۵ ماتریس کلید رمزی \lr{K1, K2, K3, K4, K5} شکل گیرند:\par

% SOURCE Hyperchaos_5D_Chapter3.txt:23-23 [spacing]


% SOURCE Hyperchaos_5D_Chapter3.txt:24-24 [paragraph]
{-} \lr{K1}\lr{(i)} = \lr{floor}\lr{( \SourceGlyph{007C}x\_i\SourceGlyph{007C} * 10\SourceGlyph{005E}14 )} \lr{mod 256}\par

% SOURCE Hyperchaos_5D_Chapter3.txt:25-25 [paragraph]
{-} \lr{K2}\lr{(i)} = \lr{floor}\lr{( \SourceGlyph{007C}y\_i\SourceGlyph{007C} * 10\SourceGlyph{005E}14 )} \lr{mod 256}\par

% SOURCE Hyperchaos_5D_Chapter3.txt:26-26 [paragraph]
{-} \lr{K3}\lr{(i)} = \lr{floor}\lr{( \SourceGlyph{007C}z\_i\SourceGlyph{007C} * 10\SourceGlyph{005E}14 )} \lr{mod 256}\par

% SOURCE Hyperchaos_5D_Chapter3.txt:27-27 [paragraph]
{-} \lr{K4}\lr{(i)} = \lr{floor}\lr{( \SourceGlyph{007C}u\_i\SourceGlyph{007C} * 10\SourceGlyph{005E}14 )} \lr{mod 256}\par

% SOURCE Hyperchaos_5D_Chapter3.txt:28-28 [paragraph]
{-} \lr{K5}\lr{(i)} = \lr{floor}\lr{( \SourceGlyph{007C}v\_i\SourceGlyph{007C} * 10\SourceGlyph{005E}14 )} \lr{mod 256}\par

% SOURCE Hyperchaos_5D_Chapter3.txt:29-29 [spacing]


% SOURCE Hyperchaos_5D_Chapter3.txt:30-30 [paragraph]
کلیدهای \lr{K1} و \lr{K2} جهت تعیین الگوی آشفته‌سازی زیگ‌زاگی و کلیدهای \lr{K3, K4, K5} جهت انتخاب پویای قواعد کدگذاری \lr{DNA} و اجرای عملگرهای انتشار \lr{DNA XOR} مورد استفاده قرار می‌گیرند. \par

% SOURCE Hyperchaos_5D_Chapter3.txt:31-31 [spacing]


% SOURCE Hyperchaos_5D_Chapter3.txt:32-32 [paragraph]
الگوریتم گام‌به‌گام راه‌اندازی سیستم ۵ بعدی و استخراج کلیدها در شبه‌کد ۳{-}۲ ارائه شده است.\par

% SOURCE Hyperchaos_5D_Chapter3.txt:33-33 [spacing]


% SOURCE Hyperchaos_5D_Chapter3.txt:34-58 [pseudocode]
\Needspace{5\baselineskip}\noindent [شبه‌کد ۳{-}۲: راه‌اندازی سیستم ابرآشوبی ۵ بعدی با حل عددی \lr{RK4} و گسسته‌سازی کلیدهای رمزی]\par\begin{sourcecode}
\begin{LTR}\ttfamily\fontsize{9}{13}\selectfont\raggedright
\CodeLine{0.0}{Input  : Initial Values (x0, y0, z0, u0, v0), Image Dimensions W, H, Step Size h = 0.001, N0 = 1000}
\CodeLine{0.0}{Output : 8{-}bit Key Matrices K1, K2, K3, K4, K5 of size W x H}
\CodeLine{0.0}{\strut}
\CodeLine{0.0}{1. Total\_\allowbreak{}Steps = N0 + (W * H)}
\CodeLine{0.0}{2. Initialize State Vector: X = [x0, y0, z0, u0, v0]\SourceGlyph{005E}T}
\CodeLine{0.0}{3. For n = 1 to Total\_\allowbreak{}Steps:}
\CodeLine{2.5}{Compute k1 = f(X)}
\CodeLine{2.5}{Compute k2 = f(X + (h/2)*k1)}
\CodeLine{2.5}{Compute k3 = f(X + (h/2)*k2)}
\CodeLine{2.5}{Compute k4 = f(X + h*k3)}
\CodeLine{2.5}{X\_\allowbreak{}next = X + (h/6) * (k1 + 2*k2 + 2*k3 + k4)}
\CodeLine{2.5}{Store X\_\allowbreak{}next to Raw\_\allowbreak{}Stream[n]}
\CodeLine{2.5}{X = X\_\allowbreak{}next}
\CodeLine{0.0}{4. Discard first N0 transient samples: Valid\_\allowbreak{}Stream = Raw\_\allowbreak{}Stream[N0+1 : Total\_\allowbreak{}Steps]}
\CodeLine{0.0}{5. For i = 1 to (W * H):}
\CodeLine{2.5}{K1[i] = floor( abs(Valid\_\allowbreak{}Stream[i].x) * 10\SourceGlyph{005E}14 ) mod 256}
\CodeLine{2.5}{K2[i] = floor( abs(Valid\_\allowbreak{}Stream[i].y) * 10\SourceGlyph{005E}14 ) mod 256}
\CodeLine{2.5}{K3[i] = floor( abs(Valid\_\allowbreak{}Stream[i].z) * 10\SourceGlyph{005E}14 ) mod 256}
\CodeLine{2.5}{K4[i] = floor( abs(Valid\_\allowbreak{}Stream[i].u) * 10\SourceGlyph{005E}14 ) mod 256}
\CodeLine{2.5}{K5[i] = floor( abs(Valid\_\allowbreak{}Stream[i].v) * 10\SourceGlyph{005E}14 ) mod 256}
\CodeLine{0.0}{6. Reshape K1, K2, K3, K4, K5 into W x H matrices}
\CodeLine{0.0}{7. Return K1, K2, K3, K4, K5}
\end{LTR}
\end{sourcecode}

% SOURCE Hyperchaos_5D_Chapter3.txt:59-59 [spacing]

% SOURCE Zigzag_DNA_Chapter3.txt:1-1 [heading]
\SourceHeading{1}{۳{-}۵. الگوریتم دو مرحله‌ای رمزنگاری (جایگشت زیگ‌زاگی و انتشار پویای \lr{DNA})}

% SOURCE Zigzag_DNA_Chapter3.txt:2-2 [spacing]


% SOURCE Zigzag_DNA_Chapter3.txt:3-3 [paragraph]
برای دستیابی به بالاترین سطح امنیت، خنثی‌سازی کامل همبستگی میان پیکسل‌های مجاور، دستیابی به آنتروپی آرمانی ($H(m) \approx 8.0$) و محقق ساختن اثر بهمنی کامل \lr{(Avalanche Effect)}، فرایند رمزنگاری تصویر اصلی طی دو مرحله متوالی و مکمل جایگشت \lr{(Confusion)} و انتشار \lr{(Diffusion)} بر پایه معماری شانون انجام می‌پذیرد \ThesisCite{1}, \ThesisCite{6}, \ThesisCite{9}.\par

% SOURCE Zigzag_DNA_Chapter3.txt:4-4 [spacing]


% SOURCE Zigzag_DNA_Chapter3.txt:5-5 [image]
\SourcePicture{chapter3_zigzag_dna_pipeline.png}{شکل ۳{-}۷: نمودار فرآیند دو مرحله‌ای جایگشت زیگ‌زاگی، جدول قواعد ۸گانه \lr{DNA} و انتشار پویای \lr{DNA}}{[شکل ۳{-}۷: نمودار فرآیند دو مرحله‌ای جایگشت زیگ‌زاگی، جدول قواعد ۸گانه \lr{DNA} و انتشار پویای \lr{DNA {-} chapter3\_zigzag\_dna\_pipeline.png}]}{0}

% SOURCE Zigzag_DNA_Chapter3.txt:6-6 [spacing]


% SOURCE Zigzag_DNA_Chapter3.txt:7-7 [heading]
\SourceHeading{2}{۳{-}۵{-}۱. مرحله اول: جایگشت زیگ‌زاگی \lr{(Zig{-}zag Scrambling / Confusion Stage)}}

% SOURCE Zigzag_DNA_Chapter3.txt:8-8 [paragraph]
در مرحله جایگشت، ماتریس تصویر ورودی واضح $I$ با ابعاد $W \times H$ با استفاده از آشفته‌سازی مکانی پویا بر اساس مرتب‌سازی اندیس‌های کلید اول سیستم ابرآشوبی ($K_1$) آشفته می‌گردد. در الگوریتم‌های سنتی رمزشده، جایگشت‌های سطری یا ستونی ساده به تنهایی توانایی شکستن همبستگی دو‌بعدی پیکسل‌ها را ندارند؛ اما استفاده از تابع \lr{argsort} روی بردار کلید شبه‌تصادفی $K_1$ برای ایجاد جایگشت یک‌تایی بر روی قطرهای فرعی و اصلی، ساختار هندسی تصویر را تماماً برهم زده و ضریب همبستگی پیکسل‌های مجاور ($r_{xy}$) را در سه جهت افقی، عمودی و قطری از حدود $1.0$ به نزدیک صفر ($r_{xy} \approx 0.0$) کاهش می‌دهد.\par

% SOURCE Zigzag_DNA_Chapter3.txt:9-9 [spacing]


% SOURCE Zigzag_DNA_Chapter3.txt:10-10 [paragraph]
الگوریتم پیمایش زیگ‌زاگی پیکسل‌های $(i,\allowbreak{} j)$ را به یک توالی یک‌بعدی آشفته‌شده $P_{\mathrm{scrambled}}$ تبدیل می‌نماید بدون آنکه شدت روشنایی و مقادیر عددی پیکسل‌ها تغییر یابند.\par

% SOURCE Zigzag_DNA_Chapter3.txt:11-11 [spacing]


% SOURCE Zigzag_DNA_Chapter3.txt:12-12 [heading]
\SourceHeading{2}{۳{-}۵{-}۲. مرحله دوم: انتشار پویای \lr{DNA} \lr{(Dynamic DNA Encoding \& XOR / Diffusion Stage)}}

% SOURCE Zigzag_DNA_Chapter3.txt:13-13 [paragraph]
در مرحله انتشار، مقادیر عددی و شدت روشنایی پیکسل‌ها تحت تغییرات غیرخطی گسترده قرار می‌گیرند به‌گونه‌ای که تغییر حتی یک پیکسل در تصویر ورودی واضح، منجر به تغییر بیش از ۵۰\SourceGlyph{066A} بیت‌های کل خروجی رمزشده گردد (محقق ساختن مقادیر آرمانی \lr{NPCR \SourceGlyph{003E} 99.60\%} و \lr{UACI} \ensuremath{\approx} \lr{33.46\%}).\par

% SOURCE Zigzag_DNA_Chapter3.txt:14-14 [spacing]


% SOURCE Zigzag_DNA_Chapter3.txt:15-15 [paragraph]
فرایند انتشار پویای \lr{DNA} شامل گام‌های زیر است:\par

% SOURCE Zigzag_DNA_Chapter3.txt:16-16 [paragraph]
۱. کدگذاری پویای \lr{DNA}: هر پیکسل ۸ بیتی تصویر آشفته‌شده شامل ۴ جفت‌بیت ۲ بیتی است. بر پایه اصول مكمل‌سازی واتسون{-}کریک ($A \leftrightarrow T$ و $C \leftrightarrow G$)، ۸ قاعده استاندارد کدگذاری \lr{DNA} تعریف می‌شوند. انتخاب قاعده کدگذاری برای هر پیکسل به صورت پویا و بر اساس توالی‌های حاصل از سیستم ۵ بعدی ابرآشوبی تعیین می‌گردد:\par

% SOURCE Zigzag_DNA_Chapter3.txt:17-17 [spacing]


% SOURCE Zigzag_DNA_Chapter3.txt:18-18 [paragraph]
\[\mathrm{Rule}_i = (S_i \bmod 8) + 1\]\par

% SOURCE Zigzag_DNA_Chapter3.txt:19-19 [spacing]


% SOURCE Zigzag_DNA_Chapter3.txt:20-20 [paragraph]
جدول قواعد ۸گانه کدگذاری \lr{DNA} در ادامه تبیین شده است:\par

% SOURCE Zigzag_DNA_Chapter3.txt:21-21 [paragraph]
{-} قاعده ۱: \lr{A=00, C=01, G=10, T=11}\par

% SOURCE Zigzag_DNA_Chapter3.txt:22-22 [paragraph]
{-} قاعده ۲: \lr{A=00, G=01, C=10, T=11}\par

% SOURCE Zigzag_DNA_Chapter3.txt:23-23 [paragraph]
{-} قاعده ۳: \lr{C=00, A=01, T=10, G=11}\par

% SOURCE Zigzag_DNA_Chapter3.txt:24-24 [paragraph]
{-} قاعده ۴: \lr{C=00, T=01, A=10, G=11}\par

% SOURCE Zigzag_DNA_Chapter3.txt:25-25 [paragraph]
{-} قاعده ۵: \lr{G=00, A=01, T=10, C=11}\par

% SOURCE Zigzag_DNA_Chapter3.txt:26-26 [paragraph]
{-} قاعده ۶: \lr{G=00, T=01, A=10, C=11}\par

% SOURCE Zigzag_DNA_Chapter3.txt:27-27 [paragraph]
{-} قاعده ۷: \lr{T=00, C=01, G=10, A=11}\par

% SOURCE Zigzag_DNA_Chapter3.txt:28-28 [paragraph]
{-} قاعده ۸: \lr{T=00, G=01, C=10, A=11}\par

% SOURCE Zigzag_DNA_Chapter3.txt:29-29 [spacing]


% SOURCE Zigzag_DNA_Chapter3.txt:30-30 [paragraph]
۲. چرخش و جهش پویای بازهای \lr{DNA} \lr{(Dynamic DNA Flip)}: بازهای نیتروژنی استخراج‌شده بر اساس مقدار متغیر حالت سیستم ۵ بعدی تحت چرخش پویای موقعیت قرار می‌گیرند.\par

% SOURCE Zigzag_DNA_Chapter3.txt:31-31 [paragraph]
۳. عملگر منطقی زیستی \lr{DNA XOR}: توالی ۴ حرفی \lr{DNA} تصویر با توالی ۴ حرفی \lr{DNA} حاصل از گسسته‌سازی سیستم ۵ بعدی ابرآشوبی تحت عملگر \lr{XOR} زیستی قرار می‌گیرند.\par

% SOURCE Zigzag_DNA_Chapter3.txt:32-32 [paragraph]
۴. رمزگشایی پویای \lr{DNA}: توالی حاصل با قاعده پویای متناظر به ماتریس تصویری رمزشده $I_{\mathrm{encrypted}}$ بازگردانده می‌شود.\par

% SOURCE Zigzag_DNA_Chapter3.txt:33-33 [spacing]


% SOURCE Zigzag_DNA_Chapter3.txt:34-34 [image]
\SourcePicture{eq3_3_zigzag_dna.png}{رابطه ۳{-}۴: فرمول‌بندی ریاضی جایگشت زیگ‌زاگی، انتخاب پویای قواعد \lr{DNA} و عملگر منطقی \lr{DNA XOR}}{[رابطه ۳{-}۴: فرمول‌بندی ریاضی جایگشت زیگ‌زاگی، انتخاب پویای قواعد \lr{DNA} و عملگر منطقی \lr{DNA XOR {-} eq3\_3\_zigzag\_dna.png}]}{2}

% SOURCE Zigzag_DNA_Chapter3.txt:35-35 [spacing]


% SOURCE Zigzag_DNA_Chapter3.txt:36-36 [paragraph]
الگوریتم کامل دو مرحله‌ای جایگشت زیگ‌زاگی و انتشار پویای \lr{DNA} در شبه‌کد ۳{-}۳ ارائه‌شده است.\par

% SOURCE Zigzag_DNA_Chapter3.txt:37-37 [spacing]


% SOURCE Zigzag_DNA_Chapter3.txt:38-64 [pseudocode]
\Needspace{5\baselineskip}\noindent [شبه‌کد ۳{-}۳: الگوریتم دو مرحله‌ای رمزنگاری شامل جایگشت زیگ‌زاگی و انتشار پویای \lr{DNA}]\par\begin{sourcecode}
\begin{LTR}\ttfamily\fontsize{9}{13}\selectfont\raggedright
\CodeLine{0.0}{Input  : Scrambled Image Matrix I\_\allowbreak{}scrambled (W x H), Chaotic Sequences K1..K5}
\CodeLine{0.0}{Output : Encrypted Image Matrix I\_\allowbreak{}encrypted (W x H)}
\CodeLine{0.0}{\strut}
\CodeLine{0.0}{1. Compute Zig{-}zag Scrambling Matrix:}
\CodeLine{1.5}{I\_\allowbreak{}scrambled = ZigZag\_\allowbreak{}Permute(I\_\allowbreak{}plain)}
\CodeLine{0.0}{\strut}
\CodeLine{0.0}{2. Convert Image Pixels to Binary Streams:}
\CodeLine{1.5}{For each pixel P(i,j) in I\_\allowbreak{}scrambled:}
\CodeLine{3.0}{Rule\_\allowbreak{}idx = (K1(i,j) mod 8) + 1}
\CodeLine{3.0}{DNA\_\allowbreak{}pixel = Encode\_\allowbreak{}DNA(P(i,j), Rule\_\allowbreak{}idx)  // Returns 4 DNA bases: e.g. "ACGT"}
\CodeLine{3.0}{\strut}
\CodeLine{3.0}{DNA\_\allowbreak{}chaotic = Encode\_\allowbreak{}DNA(K2(i,j), Rule\_\allowbreak{}idx) // Returns 4 DNA bases}
\CodeLine{3.0}{\strut}
\CodeLine{3.0}{// Dynamic DNA XOR Operation}
\CodeLine{3.0}{DNA\_\allowbreak{}cipher = DNA\_\allowbreak{}XOR\_\allowbreak{}Operation(DNA\_\allowbreak{}pixel, DNA\_\allowbreak{}chaotic)}
\CodeLine{3.0}{\strut}
\CodeLine{3.0}{// Dynamic DNA Base Flip (Mutation)}
\CodeLine{3.0}{If (K3(i,j) mod 2 == 1):}
\CodeLine{5.0}{DNA\_\allowbreak{}cipher = DNA\_\allowbreak{}Complement\_\allowbreak{}Flip(DNA\_\allowbreak{}cipher)}
\CodeLine{5.0}{\strut}
\CodeLine{3.0}{// Decode back to 8{-}bit integer}
\CodeLine{3.0}{I\_\allowbreak{}encrypted(i,j) = Decode\_\allowbreak{}DNA(DNA\_\allowbreak{}cipher, Rule\_\allowbreak{}idx)}
\CodeLine{0.0}{\strut}
\CodeLine{0.0}{3. Return Encrypted Image Matrix I\_\allowbreak{}encrypted}
\end{LTR}
\end{sourcecode}

% SOURCE Zigzag_DNA_Chapter3.txt:65-65 [spacing]

% SOURCE Saliency_Steganography_Chapter3.txt:1-1 [heading]
\SourceHeading{1}{۳{-}۶. ماژول پنهان‌نگاری فراداده مبتنی بر برجستگی بصری \lr{(Saliency{-}aware Steganography)} و فشرده‌سازی \lr{zlib}}

% SOURCE Saliency_Steganography_Chapter3.txt:2-2 [spacing]


% SOURCE Saliency_Steganography_Chapter3.txt:3-3 [paragraph]
در سامانه‌های اینترنت اشیای پزشکی \lr{(IoMT)} و پزشکی از راه دور، صیانت از حریم خصوصی بیمار و حفاظت از فراداده‌های بالینی و هویتی (شامل اطلاعات \lr{EXIF}، پرونده بالینی، نام و شناسه بیمار) همانند محرمانگی خودِ تصویر تشخیصی از اهمیت حیاتی برخوردار است. با این حال، الگوریتم‌های سنتی پنهان‌نگاری (مانند \lr{LSB} ساده یا \lr{PVD}) داده‌ها را به صورت غیرتطبیقی در سراسر تصویر جای‌دهی می‌کنند که این امر منجر به ایجاد نویز دیداری، اعوجاج آناتومیک و افت کیفیت بالینی در ناحیه حساس تشخیصی \lr{(ROI)} می‌گردد \ThesisCite{8}, \ThesisCite{9}.\par

% SOURCE Saliency_Steganography_Chapter3.txt:4-4 [spacing]


% SOURCE Saliency_Steganography_Chapter3.txt:5-5 [paragraph]
برای غلبه بر این محدودیت بنیادی و نیل به هدف حفظ کامل اصالت بافت‌های تشخیصی، سیستم پیشنهادی از یک پروتکل پنهان‌نگاری تطبیقی دو مرحله‌ای مبتنی بر تحلیل برجستگی بصری \lr{(Saliency Detection)} و فشرده‌سازی بدون اتلاف \lr{zlib} استفاده می‌نماید. روند کارهای این ماژول در شکل ۳{-}۸ به تصویر کشیده شده است.\par

% SOURCE Saliency_Steganography_Chapter3.txt:6-6 [spacing]


% SOURCE Saliency_Steganography_Chapter3.txt:7-7 [image]
\SourcePicture{chapter3_saliency_stego_pipeline.png}{شکل ۳{-}۸: نمودار فرآیند پنهان‌نگاری فراداده بیمار مبتنی بر فشرده‌سازی \lr{zlib} و برجستگی بصری}{[شکل ۳{-}۸: نمودار فرآیند پنهان‌نگاری فراداده بیمار مبتنی بر فشرده‌سازی \lr{zlib} و برجستگی بصری {-} \lr{chapter3\_saliency\_stego\_pipeline.png}]}{0}

% SOURCE Saliency_Steganography_Chapter3.txt:8-8 [spacing]


% SOURCE Saliency_Steganography_Chapter3.txt:9-9 [paragraph]
فرایند اجرایی ماژول پنهان‌نگاری فراداده طی سه گام اصلی انجام می‌پذیرد:\par

% SOURCE Saliency_Steganography_Chapter3.txt:10-10 [spacing]


% SOURCE Saliency_Steganography_Chapter3.txt:11-11 [paragraph]
۱. فشرده‌سازی بدون اتلاف \lr{zlib} \lr{(zlib Lossless Compression)}:\par

% SOURCE Saliency_Steganography_Chapter3.txt:12-12 [paragraph]
متن فراداده‌های هویتی بیمار $M_{\mathrm{patient}}$ ابتدا توسط الگوریتم استاندارد فشرده‌سازی بدون اتلاف \lr{zlib} پردازش می‌شود. این امر دو مزیت کلیدی را فراهم می‌سازد: نخست، کاهش حجم داده‌های جای‌دهی تا حدود ۶۵\SourceGlyph{066A} که ظرفیت مورد نیاز را به شدت کاهش می‌دهد؛ دوم، افزایش درجه آنتروپی و تصادفی بودن رشته بیت‌های محرمانه که پنهان‌نگاری را در برابر حملات تحلیل پنهان‌کاوی \lr{(Steganalysis)} فوق‌العاده مقاوم می‌سازد:\par

% SOURCE Saliency_Steganography_Chapter3.txt:13-13 [spacing]


% SOURCE Saliency_Steganography_Chapter3.txt:14-14 [paragraph]
\[\mathrm{Payload}_{\mathrm{compressed}} = \mathrm{zlib\_compress}(M_{\mathrm{patient}})\]\par

% SOURCE Saliency_Steganography_Chapter3.txt:15-15 [spacing]


% SOURCE Saliency_Steganography_Chapter3.txt:16-16 [paragraph]
۲. استخراج نقشه برجستگی بصری \lr{(Spectral Saliency Detection)}:\par

% SOURCE Saliency_Steganography_Chapter3.txt:17-17 [paragraph]
با استفاده از الگوریتم تحلیل طیفی برجستگی بصری \lr{(Spectral Residual Saliency)}، نقشه توجه دیداری $S(x,\allowbreak{}y)$ برای تصویر رمزشده محاسبه می‌شود. طیف برجستگی با اعمال تبدیل فوریه سریع \lr{(FFT)}، استخراج دامنه طیفی و محاسبه پس‌ماند طیفی \lr{(Spectral Residual)} در حوزه فرکانس به دست می‌آید. پیکسل‌های تصویر بر اساس آستانه دینامیکی $T_{\mathrm{saliency}}$ به دو دسته تفکیک می‌شوند:\par

% SOURCE Saliency_Steganography_Chapter3.txt:18-18 [paragraph]
{-} نواحی با برجستگی بالا \lr{(High Saliency / ROI)}: پیکسل‌هایی که حاوی اطلاعات آناتومیک و تشخیصی حساس هستند.\par

% SOURCE Saliency_Steganography_Chapter3.txt:19-19 [paragraph]
{-} نواحی کم‌برجستگی \lr{(Low Saliency / Background)}: پیکسل‌های متعلق به پس‌زمینه بی‌اهمیت تصویر پزشکی.\par

% SOURCE Saliency_Steganography_Chapter3.txt:20-20 [spacing]


% SOURCE Saliency_Steganography_Chapter3.txt:21-21 [paragraph]
۳. جای‌دهی تطبیقی در بیت‌های کم‌ارزش \lr{(Adaptive 2{-}bit LSB Embedding)}:\par

% SOURCE Saliency_Steganography_Chapter3.txt:22-22 [paragraph]
جریان بیت‌های فشرده‌سازی شده $\mathrm{Payload}_{\mathrm{compressed}}$ صرفاً در ۲ بیت کم‌ارزش \lr{(LSB)} پیکسل‌هایی جای‌دهی می‌شوند که شرط قرارگیری در ناحیه پس‌زمینه کم‌برجستگی را برآورده سازند ($M_{\mathrm{ROI}}(i,\allowbreak{}j) = 0$). این جای‌دهی دقیقاً تضمین می‌کند که ناحیه حساس \lr{ROI} ۱۰۰\SourceGlyph{066A} دست‌نخورده باقی بماند ($PSNR_{\mathrm{ROI}} = \infty$) و شاخص‌های کلی تصویر پنهان‌نگاری‌شده در سطح آرمانی $PSNR > 50\mathrm{~dB}$ و $SSIM \approx 1.0$ حفظ گردند \ThesisCite{3}, \ThesisCite{8}, \ThesisCite{9}.\par

% SOURCE Saliency_Steganography_Chapter3.txt:23-23 [spacing]


% SOURCE Saliency_Steganography_Chapter3.txt:24-24 [paragraph]
فرمول‌بندی جامع ریاضی، روابط فشرده‌سازی \lr{zlib}، استخراج طیف برجستگی و جای‌دهی \lr{LSB} در رابطه ۳{-}۵ نشان داده شده است.\par

% SOURCE Saliency_Steganography_Chapter3.txt:25-25 [spacing]


% SOURCE Saliency_Steganography_Chapter3.txt:26-26 [image]
\SourcePicture{eq3_4_saliency_stego.png}{رابطه ۳{-}۵: فرمول‌بندی نقشه برجستگی بصری، فشرده‌سازی \lr{zlib} و جای‌دهی پس‌زمینه}{[رابطه ۳{-}۵: فرمول‌بندی نقشه برجستگی بصری، فشرده‌سازی \lr{zlib} و جای‌دهی پس‌زمینه {-} \lr{eq3\_4\_saliency\_stego.png}]}{2}

% SOURCE Saliency_Steganography_Chapter3.txt:27-27 [spacing]


% SOURCE Saliency_Steganography_Chapter3.txt:28-28 [paragraph]
الگوریتم گام‌به‌گام پنهان‌نگاری فراداده بیمار آگاه به برجستگی بصری در شبه‌کد ۳{-}۴ ارائه شده است.\par

% SOURCE Saliency_Steganography_Chapter3.txt:29-29 [spacing]


% SOURCE Saliency_Steganography_Chapter3.txt:30-50 [pseudocode]
\Needspace{5\baselineskip}\noindent [شبه‌کد ۳{-}۴: پنهان‌نگاری فراداده بیمار مبتنی بر برجستگی بصری و فشرده‌سازی \lr{zlib}]\par\begin{sourcecode}
\begin{LTR}\ttfamily\fontsize{9}{13}\selectfont\raggedright
\CodeLine{0.0}{Input  : Encrypted Image I\_\allowbreak{}encrypted (W x H), Patient Metadata M\_\allowbreak{}patient}
\CodeLine{0.0}{Output : Stego{-}Encrypted Image I\_\allowbreak{}stego, Embedded Bit Count}
\CodeLine{0.0}{\strut}
\CodeLine{0.0}{1. Compress Metadata: Payload = zlib\_\allowbreak{}compress(M\_\allowbreak{}patient)}
\CodeLine{0.0}{2. Convert Payload to binary stream: Bit\_\allowbreak{}Stream = StringToBits(Payload)}
\CodeLine{0.0}{3. Compute Spectral Saliency Map: S\_\allowbreak{}map = Compute\_\allowbreak{}Spectral\_\allowbreak{}Saliency(I\_\allowbreak{}encrypted)}
\CodeLine{0.0}{4. Generate Background Mask: M\_\allowbreak{}bg = (S\_\allowbreak{}map \SourceGlyph{003C} T\_\allowbreak{}saliency) AND (M\_\allowbreak{}ROI == 0)}
\CodeLine{0.0}{5. Initialize bit index: bit\_\allowbreak{}idx = 0, total\_\allowbreak{}bits = Length(Bit\_\allowbreak{}Stream)}
\CodeLine{0.0}{6. For each pixel (i, j) in I\_\allowbreak{}encrypted where M\_\allowbreak{}bg(i,j) == 1:}
\CodeLine{3.0}{If bit\_\allowbreak{}idx \SourceGlyph{003C} total\_\allowbreak{}bits:}
\CodeLine{4.5}{Extract 2 bits: b1, b2 = Bit\_\allowbreak{}Stream[bit\_\allowbreak{}idx], Bit\_\allowbreak{}Stream[bit\_\allowbreak{}idx + 1]}
\CodeLine{4.5}{Modify 2 LSBs: I\_\allowbreak{}stego(i,j) = (I\_\allowbreak{}encrypted(i,j) AND 11111100\_\allowbreak{}2) OR (b1 \SourceGlyph{007C}\SourceGlyph{007C} b2)}
\CodeLine{4.5}{bit\_\allowbreak{}idx = bit\_\allowbreak{}idx + 2}
\CodeLine{3.0}{Else:}
\CodeLine{4.5}{I\_\allowbreak{}stego(i,j) = I\_\allowbreak{}encrypted(i,j)}
\CodeLine{0.0}{7. For pixels where M\_\allowbreak{}bg(i,j) == 0:}
\CodeLine{3.0}{I\_\allowbreak{}stego(i,j) = I\_\allowbreak{}encrypted(i,j)  // 100\% Unchanged ROI}
\CodeLine{0.0}{8. Return I\_\allowbreak{}stego, bit\_\allowbreak{}idx}
\end{LTR}
\end{sourcecode}

% SOURCE Saliency_Steganography_Chapter3.txt:51-51 [spacing]

% SOURCE Complexity_Edge_Chapter3.txt:1-1 [heading]
\SourceHeading{1}{۳{-}۷. تحلیل پیچیدگی محاسباتی، شبه‌کدهای اجرایی و ملاحظات پیاده‌سازی لبه \lr{(Edge Computing)}}

% SOURCE Complexity_Edge_Chapter3.txt:2-2 [spacing]


% SOURCE Complexity_Edge_Chapter3.txt:3-3 [paragraph]
جهت تضمین تکرارپذیری \lr{(Reproducibility)} الگوریتم‌ها و اثبات قابلیت پیاده‌سازی عملیاتی سیستم امنیتی هیبریدی پیشنهادی در بسترهای واقعی سلامت الکترونیک، این بخش به ارائه شبه‌کدهای استاندارد چهارگانه، تحلیل ریاضی پیچیدگی محاسباتی زمانی و حافظه‌ای، و ارزیابی نرخ تاخیر \lr{(Latency)} بر روی سخت‌افزارهای لبه اینترنت اشیای پزشکی \lr{(IoMT)} اختصاص یافته است \ThesisCite{1}, \ThesisCite{10}.\par

% SOURCE Complexity_Edge_Chapter3.txt:4-4 [spacing]


% SOURCE Complexity_Edge_Chapter3.txt:5-5 [heading]
\SourceHeading{2}{۳{-}۷{-}۱. شبه‌کدهای اجرایی زنجیره پردازشی فرستنده}

% SOURCE Complexity_Edge_Chapter3.txt:6-6 [paragraph]
برای شفاف‌سازی دقیق گام‌های محاسباتی، لوله‌های پردازشی ۴لایه سیستم در قالب ۴ شبه‌کد استاندارد دانشگاهی تدوین گردیده‌اند:\par

% SOURCE Complexity_Edge_Chapter3.txt:7-7 [spacing]


% SOURCE Complexity_Edge_Chapter3.txt:8-8 [paragraph]
۱. شبه‌کد ۳{-}۱: استخراج \lr{ROI} و تولید کلید پویای \lr{SHA{-}256}\par

% SOURCE Complexity_Edge_Chapter3.txt:9-9 [spacing]


% SOURCE Complexity_Edge_Chapter3.txt:10-32 [pseudocode]
\Needspace{5\baselineskip}\noindent [شبه‌کد ۳{-}۱: قطعه‌بندی \lr{ROI} با \lr{U{-}Net}، استخراج ویژگی‌های آماری و کلیدسازی پویا با \lr{SHA{-}256}]\par\begin{sourcecode}
\begin{LTR}\ttfamily\fontsize{9}{13}\selectfont\raggedright
\CodeLine{0.0}{Input  : Medical Image I (W x H), Trained U{-}Net Model}
\CodeLine{0.0}{Output : Binary Mask M\_\allowbreak{}ROI, 256{-}bit Hash Key K\_\allowbreak{}hash, 5D Initial Values (x0, y0, z0, u0, v0)}
\CodeLine{0.0}{\strut}
\CodeLine{0.0}{1. M\_\allowbreak{}ROI = U\_\allowbreak{}Net\_\allowbreak{}Inference(I, Trained\_\allowbreak{}Model)}
\CodeLine{0.0}{2. Extract ROI Pixels: P\_\allowbreak{}ROI = \{ I(i,j) \SourceGlyph{007C} M\_\allowbreak{}ROI(i,j) == 1 \}}
\CodeLine{0.0}{3. Compute Statistical Moments:}
\CodeLine{1.5}{{-} mu\_\allowbreak{}ROI   = Mean(P\_\allowbreak{}ROI)}
\CodeLine{1.5}{{-} var\_\allowbreak{}ROI  = Variance(P\_\allowbreak{}ROI)}
\CodeLine{1.5}{{-} skew\_\allowbreak{}ROI = Skewness(P\_\allowbreak{}ROI)}
\CodeLine{1.5}{{-} kurt\_\allowbreak{}ROI = Kurtosis(P\_\allowbreak{}ROI)}
\CodeLine{0.0}{4. Construct Feature Stream: F\_\allowbreak{}str = String(mu\_\allowbreak{}ROI) \SourceGlyph{007C}\SourceGlyph{007C} String(var\_\allowbreak{}ROI) \SourceGlyph{007C}\SourceGlyph{007C} String(skew\_\allowbreak{}ROI) \SourceGlyph{007C}\SourceGlyph{007C} String(kurt\_\allowbreak{}ROI) \SourceGlyph{007C}\SourceGlyph{007C} String(W) \SourceGlyph{007C}\SourceGlyph{007C} String(H)}
\CodeLine{0.0}{5. K\_\allowbreak{}hash = SHA256(F\_\allowbreak{}str)  // 256{-}bit Hash Key}
\CodeLine{0.0}{6. Split K\_\allowbreak{}hash into 5 blocks of 48 bits: K\_\allowbreak{}b1, K\_\allowbreak{}b2, K\_\allowbreak{}b3, K\_\allowbreak{}b4, K\_\allowbreak{}b5}
\CodeLine{0.0}{7. Compute Initial Values:}
\CodeLine{1.5}{{-} x0 = (K\_\allowbreak{}b1 mod 10\SourceGlyph{005E}8) / 10\SourceGlyph{005E}8}
\CodeLine{1.5}{{-} y0 = (K\_\allowbreak{}b2 mod 10\SourceGlyph{005E}8) / 10\SourceGlyph{005E}8}
\CodeLine{1.5}{{-} z0 = (K\_\allowbreak{}b3 mod 10\SourceGlyph{005E}8) / 10\SourceGlyph{005E}8}
\CodeLine{1.5}{{-} u0 = (K\_\allowbreak{}b4 mod 10\SourceGlyph{005E}8) / 10\SourceGlyph{005E}8}
\CodeLine{1.5}{{-} v0 = (K\_\allowbreak{}b5 mod 10\SourceGlyph{005E}8) / 10\SourceGlyph{005E}8}
\CodeLine{0.0}{8. Return M\_\allowbreak{}ROI, K\_\allowbreak{}hash, (x0, y0, z0, u0, v0)}
\end{LTR}
\end{sourcecode}

% SOURCE Complexity_Edge_Chapter3.txt:33-33 [spacing]


% SOURCE Complexity_Edge_Chapter3.txt:34-34 [paragraph]
۲. شبه‌کد ۳{-}۲: راه‌اندازی سیستم ۵بعدی و گسسته‌سازی کلیدها\par

% SOURCE Complexity_Edge_Chapter3.txt:35-35 [spacing]


% SOURCE Complexity_Edge_Chapter3.txt:36-53 [pseudocode]
\Needspace{5\baselineskip}\noindent [شبه‌کد ۳{-}۲: راه‌اندازی سیستم ابرآشوبی ۵ بعدی با حل عددی \lr{RK4} و گسسته‌سازی کلیدهای رمزی]\par\begin{sourcecode}
\begin{LTR}\ttfamily\fontsize{9}{13}\selectfont\raggedright
\CodeLine{0.0}{Input  : Initial Values (x0, y0, z0, u0, v0), Image Dimensions W x H, Step Size h}
\CodeLine{0.0}{Output : 8{-}bit Chaotic Key Matrices K1, K2, K3, K4, K5 of size W x H}
\CodeLine{0.0}{\strut}
\CodeLine{0.0}{1. Total\_\allowbreak{}Steps = N\_\allowbreak{}0 + W * H    // N\_\allowbreak{}0 = 1000 for Transient Removal}
\CodeLine{0.0}{2. Initialize State Vector X\_\allowbreak{}0 = [x0, y0, z0, u0, v0]\SourceGlyph{005E}T}
\CodeLine{0.0}{3. For n = 0 to Total\_\allowbreak{}Steps {-} 1:}
\CodeLine{3.0}{X\_\allowbreak{}\{n+1\} = RK4\_\allowbreak{}Integrate\_\allowbreak{}5D(X\_\allowbreak{}n, h)}
\CodeLine{0.0}{4. Discard first N\_\allowbreak{}0 states (Transient Elimination)}
\CodeLine{0.0}{5. For i = 1 to W * H:}
\CodeLine{3.0}{K1[i] = floor((abs(x[N\_\allowbreak{}0 + i]) mod 1.0) * 10\SourceGlyph{005E}8) mod 256}
\CodeLine{3.0}{K2[i] = floor((abs(y[N\_\allowbreak{}0 + i]) mod 1.0) * 10\SourceGlyph{005E}8) mod 256}
\CodeLine{3.0}{K3[i] = floor((abs(z[N\_\allowbreak{}0 + i]) mod 1.0) * 10\SourceGlyph{005E}8) mod 256}
\CodeLine{3.0}{K4[i] = floor((abs(u[N\_\allowbreak{}0 + i]) mod 1.0) * 10\SourceGlyph{005E}8) mod 256}
\CodeLine{3.0}{K5[i] = floor((abs(v[N\_\allowbreak{}0 + i]) mod 1.0) * 10\SourceGlyph{005E}8) mod 256}
\CodeLine{0.0}{6. Return K1, K2, K3, K4, K5}
\end{LTR}
\end{sourcecode}

% SOURCE Complexity_Edge_Chapter3.txt:54-54 [spacing]


% SOURCE Complexity_Edge_Chapter3.txt:55-55 [paragraph]
۳. شبه‌کد ۳{-}۳: رمزنگاری زیگ‌زاگی و انتشار پویای \lr{DNA}\par

% SOURCE Complexity_Edge_Chapter3.txt:56-56 [spacing]


% SOURCE Complexity_Edge_Chapter3.txt:57-71 [pseudocode]
\Needspace{5\baselineskip}\noindent [شبه‌کد ۳{-}۳: الگوریتم دو مرحله‌ای رمزنگاری شامل جایگشت زیگ‌زاگی و انتشار پویای \lr{DNA}]\par\begin{sourcecode}
\begin{LTR}\ttfamily\fontsize{9}{13}\selectfont\raggedright
\CodeLine{0.0}{Input  : Plain Image I (W x H), Key Matrices K1, K2, K3, K4, K5}
\CodeLine{0.0}{Output : Encrypted Image I\_\allowbreak{}encrypted}
\CodeLine{0.0}{\strut}
\CodeLine{0.0}{1. I\_\allowbreak{}scrambled = ZigZag\_\allowbreak{}Scramble\_\allowbreak{}Matrix(I)}
\CodeLine{0.0}{2. For i = 1 to W * H:}
\CodeLine{3.0}{Rule\_\allowbreak{}i = (K1[i] mod 8) + 1                 // Dynamic DNA Rule 1{-}8}
\CodeLine{3.0}{DNA\_\allowbreak{}img = Encode\_\allowbreak{}DNA(I\_\allowbreak{}scrambled[i], Rule\_\allowbreak{}i)}
\CodeLine{3.0}{DNA\_\allowbreak{}key = Encode\_\allowbreak{}DNA(K2[i], Rule\_\allowbreak{}i)}
\CodeLine{3.0}{DNA\_\allowbreak{}flipped = Dynamic\_\allowbreak{}DNA\_\allowbreak{}Flip(DNA\_\allowbreak{}img, K3[i])}
\CodeLine{3.0}{DNA\_\allowbreak{}cipher  = Dynamic\_\allowbreak{}DNA\_\allowbreak{}XOR(DNA\_\allowbreak{}flipped, DNA\_\allowbreak{}key)}
\CodeLine{3.0}{I\_\allowbreak{}encrypted[i] = Decode\_\allowbreak{}DNA(DNA\_\allowbreak{}cipher, Rule\_\allowbreak{}i)}
\CodeLine{0.0}{3. Return I\_\allowbreak{}encrypted}
\end{LTR}
\end{sourcecode}

% SOURCE Complexity_Edge_Chapter3.txt:72-72 [spacing]


% SOURCE Complexity_Edge_Chapter3.txt:73-73 [paragraph]
۴. شبه‌کد ۳{-}۴: پنهان‌نگاری \lr{Saliency} فراداده\par

% SOURCE Complexity_Edge_Chapter3.txt:74-74 [spacing]


% SOURCE Complexity_Edge_Chapter3.txt:75-85 [pseudocode]
\Needspace{5\baselineskip}\noindent [شبه‌کد ۳{-}۴: پنهان‌نگاری فراداده بیمار مبتنی بر برجستگی بصری و فشرده‌سازی \lr{zlib}]\par\begin{sourcecode}
\begin{LTR}\ttfamily\fontsize{9}{13}\selectfont\raggedright
\CodeLine{0.0}{Input  : Encrypted Image I\_\allowbreak{}encrypted, Patient Metadata M\_\allowbreak{}patient}
\CodeLine{0.0}{Output : Stego Image I\_\allowbreak{}stego}
\CodeLine{0.0}{\strut}
\CodeLine{0.0}{1. Payload\_\allowbreak{}zlib = zlib\_\allowbreak{}compress(M\_\allowbreak{}patient)}
\CodeLine{0.0}{2. S\_\allowbreak{}map = Compute\_\allowbreak{}Spectral\_\allowbreak{}Residual\_\allowbreak{}Saliency(I\_\allowbreak{}encrypted)}
\CodeLine{0.0}{3. Mask\_\allowbreak{}bg = Threshold\_\allowbreak{}Low\_\allowbreak{}Saliency\_\allowbreak{}Pixels(S\_\allowbreak{}map)}
\CodeLine{0.0}{4. I\_\allowbreak{}stego = Embed\_\allowbreak{}2bit\_\allowbreak{}LSB\_\allowbreak{}In\_\allowbreak{}Background(I\_\allowbreak{}encrypted, Payload\_\allowbreak{}zlib, Mask\_\allowbreak{}bg)}
\CodeLine{0.0}{5. Return I\_\allowbreak{}stego}
\end{LTR}
\end{sourcecode}

% SOURCE Complexity_Edge_Chapter3.txt:86-86 [spacing]


% SOURCE Complexity_Edge_Chapter3.txt:87-87 [heading]
\SourceHeading{2}{۳{-}۷{-}۲. تحلیل ریاضی پیچیدگی محاسباتی \lr{(Computational Complexity)}}

% SOURCE Complexity_Edge_Chapter3.txt:88-88 [paragraph]
برای ارزیابی کارایی عملیاتی سیستم، پیچیدگی زمانی و حافظه‌ای لایه‌های چهارگانه با نماد \lr{Big{-}O} فرمول‌بندی شده است:\par

% SOURCE Complexity_Edge_Chapter3.txt:89-89 [spacing]


% SOURCE Complexity_Edge_Chapter3.txt:90-90 [paragraph]
{-} \textbf{پیچیدگی زمانی لایه اول \lr{(U{-}Net \& SHA{-}256)}:} استنتاج شبکه \lr{U{-}Net} بر روی ماتریس $W \times H$ دارای پیچیدگی $O(W \times H)$ است. استخراج ۴ گشتاور آماری و محاسبه چکیده \lr{SHA{-}256} بر روی $N_{\mathrm{ROI}}$ پیکسل انجام شده که $N_{\mathrm{ROI}} < W \times H$ بوده و پیچیدگی $O(N_{\mathrm{ROI}})$ را تحمیل می‌نماید.\par

% SOURCE Complexity_Edge_Chapter3.txt:91-91 [paragraph]
{-} \textbf{پیچیدگی زمانی لایه دوم (سیستم ۵بعدی \lr{RK4}):} حل عددی سیستم ۵بعدی برای $N = W \times H$ گام دارای پیچیدگی خطی $O(W \times H)$ است.\par

% SOURCE Complexity_Edge_Chapter3.txt:92-92 [paragraph]
{-} \textbf{پیچیدگی زمانی لایه سوم (زیگ‌زاگی و \lr{DNA}):} پیمایش ماتریسی زیگ‌زاگی و اعمال عملگرهای \lr{DNA} برای $W \times H$ پیکسل دارای پیچیدگی خطی $O(W \times H)$ می‌باشد.\par

% SOURCE Complexity_Edge_Chapter3.txt:93-93 [paragraph]
{-} \textbf{پیچیدگی زمانی لایه چهارم \lr{(Saliency \& zlib)}:} محاسبه طیف پس‌ماند برجستگی با تبدیل فوریه سریع \lr{(FFT)} دارای پیچیدگی $O(W \times H \log(W \times H))$ و جای‌دهی \lr{LSB} دارای پیچیدگی $O(W \times H)$ است.\par

% SOURCE Complexity_Edge_Chapter3.txt:94-94 [spacing]


% SOURCE Complexity_Edge_Chapter3.txt:95-95 [paragraph]
بنابراین، پیچیدگی زمانی کل سیستم از مرتبه شبه‌خطی $O(W \times H \log(W \times H))$ بوده و پیچیدگی حافظه‌ای آن $O(W \times H)$ می‌باشد.\par

% SOURCE Complexity_Edge_Chapter3.txt:96-96 [spacing]


% SOURCE Complexity_Edge_Chapter3.txt:97-97 [image]
\SourcePicture{eq3_5_complexity_math.png}{رابطه ۳{-}۶: فرمول‌بندی ریاضی پیچیدگی محاسباتی و بودجه زمانی لایه‌ها}{[رابطه ۳{-}۶: فرمول‌بندی ریاضی پیچیدگی محاسباتی و بودجه زمانی لایه‌ها {-} \lr{eq3\_5\_complexity\_math.png}]}{2}

% SOURCE Complexity_Edge_Chapter3.txt:98-98 [spacing]


% SOURCE Complexity_Edge_Chapter3.txt:99-99 [heading]
\SourceHeading{2}{۳{-}۷{-}۳. ملاحظات پیاده‌سازی سخت‌افزاری بر روی تراشه‌های لبه \lr{IoMT} \lr{(NVIDIA Jetson TX2)}}

% SOURCE Complexity_Edge_Chapter3.txt:100-100 [paragraph]
جهت به‌کارگیری سیستم در بسترهای واقعی اینترنت اشیای پزشکی \lr{(IoMT)}، کلیه ماژول‌ها بر روی گره لبه \lr{NVIDIA Jetson TX2} (شامل پردازنده ۶ هسته‌ای \lr{ARM Cortex{-}A57}، پردازنده گرافیکی ۲۵۶ هسته‌ای با معماری \lr{Pascal} و ۴ گیگابایت حافظه \lr{LPDDR4} با توان مصرفی زیر ۱۵ وات) شبیه‌سازی و ارزیابی گردیدند \ThesisCite{10}.\par

% SOURCE Complexity_Edge_Chapter3.txt:101-101 [spacing]


% SOURCE Complexity_Edge_Chapter3.txt:102-102 [paragraph]
نتایج ارزیابی کمّی نشان می‌دهد که زمان اجرای تجمعی سیستم برای تصاویر استاندارد $256 \times 256$ برابر $2.93$ ثانیه است که تفکیک زمانی لایه‌ها به شرح زیر می‌باشد:\par

% SOURCE Complexity_Edge_Chapter3.txt:103-103 [paragraph]
{-} استنتاج \lr{U{-}Net} بر روی پردازنده گرافیکی لبه: $1.18$ ثانیه.\par

% SOURCE Complexity_Edge_Chapter3.txt:104-104 [paragraph]
{-} حل عددی سیستم ۵بعدی و تولید کلیدها: $0.42$ ثانیه.\par

% SOURCE Complexity_Edge_Chapter3.txt:105-105 [paragraph]
{-} آشفته‌سازی زیگ‌زاگی و انتشار \lr{DNA}: $0.83$ ثانیه.\par

% SOURCE Complexity_Edge_Chapter3.txt:106-106 [paragraph]
{-} پنهان‌نگاری \lr{Saliency} و فشرده‌سازی \lr{zlib}: $0.50$ ثانیه.\par

% SOURCE Complexity_Edge_Chapter3.txt:107-107 [spacing]


% SOURCE Complexity_Edge_Chapter3.txt:108-108 [paragraph]
دست‌یابی به زمان اجرای زیر ۳ ثانیه ($2.93\text{ s} < 3.0\text{ s}$) اثبات می‌کند که چارچوب پیشنهادی الزامات نرخ تاخیر زمان‌واقعی \lr{(Real{-}time Latency)} را در بسترهای بیمارستانی و پزشکی از راه دور کاملاً برآورده می‌سازد.\par

% SOURCE Complexity_Edge_Chapter3.txt:109-109 [spacing]


% SOURCE Complexity_Edge_Chapter3.txt:110-110 [image]
\SourcePicture{chapter3_edge_complexity_breakdown.png}{شکل ۳{-}۹: نمودار تفکیک زمان اجرای لایه‌ها و مقایسه کارایی بر روی سخت‌افزار لبه \lr{NVIDIA Jetson TX2}}{[شکل ۳{-}۹: نمودار تفکیک زمان اجرای لایه‌ها و مقایسه کارایی بر روی سخت‌افزار لبه \lr{NVIDIA Jetson TX2 {-} chapter3\_edge\_complexity\_breakdown.png}]}{0}

% SOURCE Complexity_Edge_Chapter3.txt:111-111 [spacing]

